\documentclass{article}
\usepackage[utf8]{inputenc}
\usepackage[russian]{babel}
\usepackage{hyperref}
\usepackage{amssymb}
\usepackage[a4paper]{geometry}

\usepackage[dvipsnames]{xcolor}
\begin{document}

{
    \thispagestyle{empty}
    \vspace*{\fill}
    \centering

    \Large Сборник заданий для практических занятий \\
    по курсу \\
    <<Проектирование архитектуры мобильных и веб-приложений>>\\
    \large
    \vspace{40pt}
    \vspace*{\fill}
    \newpage
}

\tableofcontents
\newpage

\section{Общие сведения}

Сборник содержит задания для практических занятий по курсу <<Проектирование архитектуры мобильных и веб-приложений>>.

Лекции проходят в форме мастер-классов, соответственно, практические занятия состоят в реализации
заданий, которые по структуре повторяют работы, разобранные на мастер-классах.


\section{Практические задания (осенний семестр)}
\subsection{Задание №1 (простейший REST API-сервис)}

Реализуйте с использованием FastAPI REST API-сервис, предоставляющий возможность добавлять,
удалять, изменять и получать все объекты согласно варианту, а также реализуйте возможность
получения части объектов для реализации поэтапной загрузки.

Данные должны храниться в SQLite (или в СУБД по вашему выбору),
обращение к данным должно осуществляться
посредством ORM-библиотеки. Осуществляйте полную валидацию корректности передаваемых данных.

\begin{enumerate}
	\item База данных студентов группы. Поля: фамилия, имя, отчество, пол, возраст.
	\item База данных расходов семьи. Поля: товар, стоимость, количество, дата.
	\item База данных загрузки аудиторий. Поля: дата и время начала, дата и время окончания, аудитория, преподаватель.
	\item База данных учета доходов и расходов предпринимателя. Поля: дата, тип операции (доход/расход), объем операции, описание,
корреспондент.
	\item База данных велоклуба. Поля: ФИО, тип велосипеда (MTB и др.), стаж участия в велоклубе.
	\item База данных рейсов авиакомпании. Поля: дата и время вылета, аэропорт вылета, аэропорт прилета, дата и время прилета,
марка самолета.
	\item База данных автобусных маршрутов. Поля: номер маршрута, номер парка, время начала и окончания движения,
длина маршрута в км.
	\item База данных электричек. Поля: вокзал, номер поезда, количество вагонов, тип (экспресс/обычный/спутник).
	\item База данных товаров интернет-магазина. Поля: название товара, категория, цена товара, описание товара.
	\item База заказов интернет-магазина. Поля: ФИО заказчика, стоимость заказа, скидка (в процентах), адрес доставки.
	\item База данных выборов. Поля: участок, кандидат, количество голосов.
	\item База данных практических работ. Поля: практическая работа, студент, номер варианта, номер уровня,
дата сдачи, оценка.
	\item База данных операторов и телеканалов. Поля: название, тип (спутник, кабель, интернет), охват (количество миллионов домохозяйств), минимальная
стоимость подписки.
	\item База данных тарифных планов оператора. Поля: название, тип вещания (обычный/HD), флаг общедоступности.
	\item База данных незаконно огороженных берегов. Поля: водный объект, регион, GPS-координаты, длина недоступного участка берега, дата фиксации нарушения.
	\item База данных временного прекращения движения в метро. Поля: дата и время начала прекращения
движения, дата и время окончания прекращения движения, станция начала, станция окончания.
	\item База данных проката фильмов. Поля: дата, время, кинотеатр, фильм, номер зала, тип сеанса (3D/IMAX/обычный).
	\item База данных эвакуированных автомобилей. Поля: улица, автостоянка, GPS-координаты,
тип нарушения (стоянка на проезжей части в месте запрета, стоянка на тротуаре, стоянка на газоне),
номер автомобиля, тип автомобиля (легковой/грузовой малой тоннажности/грузовой большой
тоннажности).
	\item База данных средних специальных учебных учреждений. Поля: название, адрес, тип подчинения (федеральный/региональный),
год основания, номер лицензии, номер аккредитации, дата окончания действия аккредитации.
	\item База данных поселков. Поля: название, девелопер, площадь, количество жителей.
	\item База данных сухопутной военной техники. Поля: название, модель, разработчик, предприятие, стоимость, тип.
	\item База данных деревьев в городе. Поля: GPS-координаты, вид дерева, округ, год посадки.
	\item База данных футбольных матчей. Поля: дата, команда, команда, счет, место проведения.
	\item База данных обращений жителей. Поля: дата, время, объект, заявитель, содержание обращения (до 255 символов),
дата ответа, ответ на обращение (до 255 символов).
	\item База данных студентов колледжа. Поля: ФИО, группа, признак бюджетности,
стипендия (нет/обычная/повышенная), флаг наличия социальной стипендии, дата рождения.
\end{enumerate}

\subsection{Задание №2 (простейший проект с использованием React)}

Реализуйте клиент с использованием React для отображения и удаления данных из сервиса,
созданного в предыдущем задании.

Стек: React, Vite, pnpm, TypeScript, Axios, Material UI.

\subsection{Задание №3 (завершение заданий №№1--2)}

Реализуйте в клиенте с использованием React:
\begin{itemize}
  \item Добавление объекта
  \item Изменение объекта
%  \item Корректную отмену операций Axios
  \item Бесконечный скроллинг
  \item Роутинг
\end{itemize}

\subsection{Задание №4}

В заданиях №№1--3 реализуйте аутентификацию, регистрацию и авторизацию пользователей (JWT).
Объекты должны отображаться, изменяться и удаляться только пользователем, который их создал.

\subsection{Задание №5}

Реализуйте сервис, полученный в заданиях №№1--4, в микросервисной
архитектуре (отдельно микросервис для аутентификации и регистрации; отдельно микросервис
для REST API заметок; отдельно микросервис для хостинга React-приложения; данные следует хранить в отдельных
контейнерах PostgreSQL).

\subsection{Задание №6}

Реализуйте автоматическое тестирование разработанных микросервисов.

\subsection{Задание №7}

Реализуйте автоматическое UI-тестирование клиентской части приложения.


\subsection{Экзамен}

Экзамен состоит в защите проекта (сервер + клиент), созданного в течение семестра.
В качестве альтернативы можно защитить свой проект, реализованный вне пар и для
других целей, в случае, если в нем реализован REST API-сервис (на любом языке)
и реализован веб-клиент с использованием реактивных фреймворков, а также используется
автоматическое тестирование (в том числе UI).


\section{Практические задания (весенний семестр)}


\subsection{Задание №1}

Реализуйте запуск Kafka в отдельном контейнере.
При любых изменениях сущности (create/update/delete) публикуйте событие в Kafka-топик
(например, ExpenseCreated, OrderPlaced).
Реализуйте простой consumer, который логирует события, а также consumer для записи в БД.
Обеспечьте идемпотентность consumers Kafka.

\subsection{Задание №2}

Реализуйте Event Sourcing для сущности (например, Expense).
Храните события в Kafka-топике.
Реализуйте replay для получения текущего состояния.
Реализуйте чистую реализацию Event Sourcing (в учебных целях), без какой-либо
оптимизации чтения.


\subsection{Задание №3}

Придумайте и реализуйте сложную бизнес-операцию в своем домене, требующую координации (например, «Оформление заказа с проверкой наличия», «Регистрация рейса с резервированием слотов», «Обработка обращения с уведомлением департамента»).
Реализуйте Saga Choreography.


\subsection{Задание №4}

\begin{enumerate}
	\item
Реализуйте Outbox pattern для всех случаев, где осуществляется запись в базу данных,
сопровождаемая записью в Kafka (кроме choreography).
	\item
Реализуйте Schema Registry для используемых при передаче информации через
Kafka данных (там, где не используется Outbox pattern). 

\textbf{Замечание}: учебное задание не подразумевает строгое соответствие стандартным 
практикам, а ориентировано на знакомство с различными подходами. Потому требование 
реализации Schema Registry только там, где нет использования Outbox pattern, просто позволяет уменьшить
объем работы.

	\item
Реализуйте механизм DLQ: при возникновении неустранимой ошибки обработки события (например, после 3 неудачных попыток)
событие должно перенаправляться в отдельный DLQ-топик для последующей ручной обработки.
\end{enumerate}


%Добавить Outbox pattern для надёжности.
%Обеспечить idempotency consumers.

%Выделить отдельный Query-микросервис.
%Реализовать проекцию всех событий в read-модель (таблицы оптимизированные под частые запросы домена).
%Перенести все GET-эндпоинты в Query-сервис.
%Обеспечить eventual consistency.
%пш

\subsection{Зачёт}

Зачёт состоит в простом подсчёте средней оценки по выполненным заданиям (учитываются срок сдачи,
качество защиты, полнота выполнения). Вместо стандартных заданий можно защищать сторонние проекты с
аналогичными архитектурными решениями.

\end{document}

